\section*{अध्याय \ref{ch:b1:periodicity} --- आवर्तिता}

\begin{solution}{exo:b1:periodicity:1}
\ce{Ca} $[\ce{Ar}]\,4\text{s}^2$: \omterm{def:b1:periodicity:table}{आवर्त} 4, \omterm{def:b1:periodicity:table}{समूह} 2, s \omterm{def:b1:periodicity:table}{ब्लॉक}। \ce{Fe}
$[\ce{Ar}]\,3\text{d}^6 4\text{s}^2$: \omterm{def:b1:periodicity:table}{आवर्त} 4, \omterm{def:b1:periodicity:table}{समूह} 8, d \omterm{def:b1:periodicity:table}{ब्लॉक}। \ce{Se}
$[\ce{Ar}]\,3\text{d}^{10} 4\text{s}^2 4\text{p}^4$: \omterm{def:b1:periodicity:table}{आवर्त} 4, \omterm{def:b1:periodicity:table}{समूह} 16, p
\omterm{def:b1:periodicity:table}{ब्लॉक}। \ce{I} $[\ce{Kr}]\,4\text{d}^{10} 5\text{s}^2 5\text{p}^5$: \omterm{def:b1:periodicity:table}{आवर्त} 5,
\omterm{def:b1:periodicity:table}{समूह} 17, p \omterm{def:b1:periodicity:table}{ब्लॉक}।
\end{solution}

\begin{solution}{exo:b1:periodicity:2}
\ce{Na}, \ce{Cl} से बड़ा है: वही \omterm{def:b1:periodicity:table}{आवर्त}, छोटा $Z^{*}$। \ce{Na},
\ce{Na+} से बड़ा है, जो अपना 3s कोश खो चुका है। \ce{Cl-}
(\qty{181}{pm}), \ce{F-} (\qty{133}{pm}) से बड़ा है: एक कोश अधिक।
\ce{Na+} (\qty{102}{pm}), \ce{Mg^{2+}} (\qty{72}{pm}) से बड़ा है: वही
दस इलेक्ट्रॉन, नाभिकीय आवेश 11 बनाम 12।
% ledger: rion:Cl-VI, rion:F-VI, rion:Na+VI, rion:Mg2+VI
\end{solution}

\begin{solution}{exo:b1:periodicity:3}
$\ce{Na} (5.14) < \ce{Al} (5.99) < \ce{Mg} (7.65) < \ce{S} (10.36) <
\ce{P} (10.49) < \ce{Ar} (15.76)$। \omterm{def:b1:periodicity:table}{आवर्त} में सामान्य वृद्धि
दो बार टूटती है: \ce{Al}, \ce{Mg} से नीचे (हटाया जाने वाला इलेक्ट्रॉन पहला
3p इलेक्ट्रॉन है, जिसकी ऊर्जा 3s से अधिक है) और \ce{S}, \ce{P} से नीचे
(हटाया जाने वाला इलेक्ट्रॉन पहला युग्मित 3p इलेक्ट्रॉन है)।
\end{solution}

\begin{solution}{exo:b1:periodicity:4}
\omterm{def:b1:periodicity:polarisability}{ध्रुवणीयता} यह मापती है कि विद्युत क्षेत्र इलेक्ट्रॉन मेघ को कितनी आसानी से
विकृत करता है। \ce{I-}, \ce{F-} से अधिक ध्रुवणीय है: वह बड़ा है
और उसके बाहरी इलेक्ट्रॉन नाभिक से अधिक दूर हैं। \ce{Cl-} ($Z = 17$ द्वारा
थामे गए 18 इलेक्ट्रॉन) \ce{Na+} ($Z = 11$ द्वारा
थामे गए 10 इलेक्ट्रॉन) से कहीं अधिक ध्रुवणीय है, जैसे ऋणायन धनायनों से होते हैं।
\end{solution}

\begin{solution}{exo:b1:periodicity:5}
अनुपात: $18.83/5.99 = 3.1$, $28.45/18.83 = 1.5$, $119.99/28.45 = 4.2$।
उछाल तीसरे इलेक्ट्रॉन के बाद आता है: तीन \omterm{def:b1:quantum-numbers:configuration}{संयोजकता इलेक्ट्रॉन}, \omterm{def:b1:periodicity:table}{समूह} 13,
\omterm{def:b1:periodicity:table}{आवर्त} 3: ऐलुमिनियम। यह \ce{Al^{3+}} बनाता है, $1\text{s}^2 2\text{s}^2
2\text{p}^6$।
% ledger: ie:Al, ie2:Al, ie3:Al, ie4:Al
\end{solution}

\begin{solution}{exo:b1:periodicity:6}
\ce{F-} का जुड़ा इलेक्ट्रॉन छोटे 2p कोश में जाता है, जहाँ वहाँ पहले से मौजूद
सात \omterm{def:b1:quantum-numbers:configuration}{संयोजकता इलेक्ट्रॉन} उसे प्रबलता से प्रतिकर्षित करते हैं; क्लोरीन के
बड़े 3p कोश में प्रतिकर्षण कम है, इसलिए अधिक ऊर्जा
मुक्त होती है। नाइट्रोजन ($2\text{p}^3$, \omorbs{u,u,u}) में अतिरिक्त इलेक्ट्रॉन
को आधे भरे \omterm{def:b1:quantum-numbers:orbital}{उपकोश} में युग्म बनाना पड़ता; मैग्नीशियम ($3\text{s}^2$) में
उसे ऊँचा 3p \omterm{def:b1:quantum-numbers:orbital}{उपकोश} शुरू करना पड़ता। किसी भी स्थिति में
ऋणायन, परमाणु और मुक्त इलेक्ट्रॉन से अधिक स्थायी नहीं है।
\end{solution}

\begin{solution}{exo:b1:periodicity:7}
$\chi_M(\ce{Na}) = (5.14 + 0.55)/2 = \qty{2.85}{eV}$; $\chi_M(\ce{Cl}) =
(12.97 + 3.61)/2 = \qty{8.29}{eV}$। अंतर बड़ा है: \ce{Na-Cl} आबंध का
इलेक्ट्रॉन युग्म लगभग पूरी तरह क्लोरीन का है, और सोडियम
क्लोराइड आयनिक है, \ce{Na+ Cl-}।
\end{solution}

\begin{solution}{exo:b1:periodicity:8}
\ce{H-F}: $\Delta\chi = 1.78$, आयनिक क्षेत्र की सीमा पर (व्यवहार में
अत्यधिक ध्रुवीय सहसंयोजी आबंध), \ce{F^{\delta-}}। \ce{C-H}: 0.35,
लगभग अध्रुवीय। \ce{Na-Cl}: 2.23, आयनिक, \ce{Cl^{-}}। \ce{C-Cl}: 0.61,
ध्रुवीय सहसंयोजी, \ce{C^{\delta+}-Cl^{\delta-}}। \ce{O-H}: 1.24, ध्रुवीय
सहसंयोजी, \ce{O^{\delta-}-H^{\delta+}}।
\end{solution}

\begin{solution}{exo:b1:periodicity:9}
समूह में नीचे जाने पर \omterm{def:b1:quantum-numbers:configuration}{संयोजकता इलेक्ट्रॉन} बड़े $n$ वाले कोश में होता है, नाभिक से
अधिक दूर और अधिक \omterm{def:b1:quantum-numbers:configuration}{क्रोड इलेक्ट्रॉनों} के \omterm{def:b1:periodicity:screening}{आवरण} में; वह अधिक
आसानी से हटता है। रुबिडियम, जिसमें एक कोश और है, की \omterm{def:b1:periodicity:ionisation-energy}{प्रथम आयनन ऊर्जा}
\qty{4.34}{eV} से कम है।
\end{solution}

\begin{solution}{exo:b1:periodicity:10}
ज़िंक $[\ce{Ar}]\,3\text{d}^{10} 4\text{s}^2$ है: उसका इलेक्ट्रॉन भरे हुए
4s से आता है; गैलियम, $\dots 4\text{s}^2 4\text{p}^1$, अपना अकेला 4p
इलेक्ट्रॉन खोता है, जिसकी ऊर्जा अधिक है और जिसका \omterm{def:b1:periodicity:screening}{आवरण} भरे 4s और 3d करते हैं। आर्सेनिक,
$4\text{p}^3$, आधे भरे \omterm{def:b1:quantum-numbers:orbital}{उपकोश} से इलेक्ट्रॉन खोता है; सेलेनियम,
$4\text{p}^4$, युग्मित इलेक्ट्रॉन खोता है, जिसे उसके साथी का प्रतिकर्षण ऊपर धकेलता
है: इसकी आयनन ऊर्जा थोड़ी कम है।
\end{solution}

\begin{solution}{exo:b1:periodicity:11}
$\Delta E = (3.98 - 2.20)^2 = 1.78^2 = \qty{3.17}{eV} = 3.17 \times 96.5 =
\qty{306}{kJ/mol}$। तब $D(\ce{H-F}) = \tfrac12(436 + 158) + 306 = 297 +
306 \approx \qty{603}{kJ/mol}$। यह बड़ा आधिक्य आंशिक आवेशों
$\ce{H^{\delta+}}$ और $\ce{F^{\delta-}}$ के बीच का अतिरिक्त आकर्षण है, जो
\omterm{def:b1:periodicity:electronegativity}{विद्युत-ऋणात्मकता} के बड़े अंतर के कारण है।
\end{solution}

\begin{solution}{exo:b1:periodicity:12}
पॉलिंग का पैमाना आबंधों A--B की ऊर्जाओं से परिभाषित है: जो तत्व
कोई आबंध नहीं बनाता उसका कोई पॉलिंग मान नहीं होता। हीलियम, नियॉन और आर्गन कोई
स्थायी यौगिक नहीं बनाते। क्रिप्टॉन और विशेषकर ज़ेनॉन के संयोजकता कोश बड़े और अधिक
ध्रुवणीय हैं और उनकी आयनन ऊर्जाएँ कम हैं
(क्रिप्टॉन के लिए \qty{14.0}{eV}, आर्गन के \qty{15.76}{eV} के मुक़ाबले); वे
फ़्लुओरीन और ऑक्सीजन द्वारा ऑक्सीकृत होते हैं और \ce{XeF2} जैसे यौगिक बनाते हैं, इसलिए
उनकी आबंध ऊर्जाएँ मापी जा सकती हैं।
% ledger: ie:Kr, ie:Ar
\end{solution}

\begin{solution}{pb:b1:periodicity:1}
\textbf{1.} हर इलेक्ट्रॉन ऐसे आयन से हटाया जाता है जो पिछले आयन से अधिक धनात्मक और
छोटा है, और जो अपने इलेक्ट्रॉनों को अधिक प्रबलता से थामता है।
\textbf{2.} $15.04/7.65 = 1.97$; $80.14/15.04 = 5.33$; $109.27/80.14 = 1.36$।
\textbf{3.} उछाल दो इलेक्ट्रॉनों के बाद आता है: X के दो \omterm{def:b1:quantum-numbers:configuration}{संयोजकता इलेक्ट्रॉन} हैं,
\omterm{def:b1:periodicity:table}{समूह} 2; \omterm{def:b1:periodicity:table}{आवर्त} 3 में यह मैग्नीशियम है।
\textbf{4.} \ce{Mg^{2+}}, $1\text{s}^2 2\text{s}^2 2\text{p}^6$, नियॉन का
विन्यास।
\textbf{5.} धातु: \omterm{def:b1:periodicity:table}{समूह} 2, सारणी के बाईं ओर, कम \omterm{def:b1:periodicity:ionisation-energy}{प्रथम
आयनन ऊर्जा} के साथ।
\textbf{6.} तीसरा इलेक्ट्रॉन 2p क्रोड से आता है, जिसका $n$ छोटा है,
जो नाभिक के बहुत पास है और जिसका \omterm{def:b1:periodicity:screening}{आवरण} लगभग नहीं होता, और वह एक
द्विधनात्मक आयन से हटाया जाता है।
% ledger: ie:Mg, ie2:Mg, ie3:Mg, ie4:Mg
\textbf{7.} सभी में दस इलेक्ट्रॉन हैं: $1\text{s}^2 2\text{s}^2 2\text{p}^6$।
\textbf{8.} वही दस इलेक्ट्रॉन नाभिकीय आवेशों 8, 9, 11,
12, 13 द्वारा थामे जाते हैं: $Z$ बढ़ने पर मेघ सिकुड़ता है।
\textbf{9.} \ce{O^{2-}}: सबसे बड़ा, क्योंकि अपने दस इलेक्ट्रॉनों के लिए उसका
नाभिकीय आवेश सबसे कम है, और वह ऋणायन है।
\textbf{10.} $140/54 = 2.6$।
\textbf{11.} $198.8/2 = \qty{99.4}{pm}$; ऋणायन \omterm{def:b1:periodicity:radius}{सहसंयोजी त्रिज्या} से $181/99.4 = 1.8$
गुना बड़ा है।
% ledger: re:Cl2, rion:Cl-VI
\textbf{12.} $\ce{Mg^{2+}} < \ce{Na+} < \ce{Cl-}$: दोनों धनायन
समइलेक्ट्रॉनी हैं (आवेश 12 बनाम 11), और \ce{Cl-} में तीसरा कोश है।
\textbf{13.} $\chi_M$ = 10.41 (\ce{F}), 8.29 (\ce{Cl}), 7.59 (\ce{Br}),
7.53 (\ce{O}), 6.26 (\ce{C}), eV में।
\textbf{14.} मुलिकेन: $\ce{F} > \ce{Cl} > \ce{Br} > \ce{O} > \ce{C}$;
पॉलिंग: $\ce{F} > \ce{O} > \ce{Cl} > \ce{Br} > \ce{C}$। ऑक्सीजन चौथे
स्थान से दूसरे पर आ जाता है।
\textbf{15.} $3.98/10.41 = 0.382$, $3.16/8.29 = 0.381$, $2.96/7.59 =
0.390$: लगभग स्थिर, क़रीब 0.38; हैलोजनों के लिए दोनों पैमाने
समानुपाती हैं।
\textbf{16.} ऑक्सीजन की \omterm{def:b1:periodicity:electron-affinity}{इलेक्ट्रॉन बंधुता} कम (\qty{1.44}{eV}) है
क्योंकि जुड़ने वाले इलेक्ट्रॉन को सघन 2p कोश में युग्म बनाना पड़ता है; मुक्त परमाणु
का मुलिकेन मान उस खिंचाव को कम आँकता है जो ऑक्सीजन अपने
आबंधों में लगाता है, और जिसे आबंधों से बना \omterm{def:b1:periodicity:electronegativity}{पॉलिंग पैमाना} सीधे मापता है।
\textbf{17.} \ce{HCl} में क्लोरीन; \ce{ClF} में फ़्लुओरीन।
\textbf{18.} $\Delta\chi = 0.96$: ध्रुवीय सहसंयोजी।
\textbf{19.} $D(\ce{H-H}) = 2 \times 218.0 = \qty{436.0}{kJ/mol}$।
\textbf{20.} $D(\ce{Cl-Cl}) = 2 \times 121.3 = \qty{242.6}{kJ/mol}$।
\textbf{21.} $D(\ce{H-Cl}) = 218.0 + 121.3 - (-92.3) = \qty{431.6}{kJ/mol}$।
\textbf{22.} $\Delta E = 431.6 - \tfrac12(436.0 + 242.6) = 431.6 - 339.3 =
\qty{92.3}{kJ/mol}$। प्रतीकों में, $D(\ce{H-Cl}) = \Delta_fH(\ce{H}) +
\Delta_fH(\ce{Cl}) - \Delta_fH(\ce{HCl})$, जबकि सममित आबंधों का
माध्य $\Delta_fH(\ce{H}) + \Delta_fH(\ce{Cl})$ है, इसलिए $\Delta E =
-\Delta_fH(\ce{HCl})$।
\textbf{23.} $92.3/96.485 = \qty{0.957}{eV}$।
\textbf{24.} $|\chi(\ce{Cl}) - \chi(\ce{H})| = \sqrt{0.957} =
\textbf{0.98}$, सारणियों के $3.16 - 2.20 = 0.96$ के मुक़ाबले: पॉलिंग का
अंकगणित, आधुनिक आँकड़ों से करने पर, सारणीबद्ध अंतर
0.02 के भीतर देता है।
\end{solution}
